\documentclass[12pt,a4paper]{report}

% --- Packages ---
\usepackage[utf8]{inputenc}
\usepackage[vietnam]{babel}
\usepackage[T1]{fontenc}
\usepackage{amsmath, amssymb}
\usepackage{graphicx}
\usepackage{hyperref}
\usepackage{listings}
\usepackage{color}
\usepackage{geometry}
\usepackage{titlesec}
\usepackage{booktabs}
\usepackage{fancyhdr}

% --- Cấu hình trang ---
\geometry{a4paper, margin=2.5cm}
\hypersetup{
    colorlinks=true,
    linkcolor=blue,
    filecolor=magenta,      
    urlcolor=cyan,
    pdftitle={Báo cáo Dự án Chatbot RAG HUS 2025},
}

% --- Định dạng Code (Listings) ---
\definecolor{codegreen}{rgb}{0,0.6,0}
\definecolor{codegray}{rgb}{0.5,0.5,0.5}
\definecolor{codepurple}{rgb}{0.58,0,0.82}
\definecolor{backcolour}{rgb}{0.95,0.95,0.92}

\lstdefinestyle{mystyle}{
    backgroundcolor=\color{backcolour},   
    commentstyle=\color{codegreen},
    keywordstyle=\color{magenta},
    numberstyle=\tiny\color{codegray},
    stringstyle=\color{codepurple},
    basicstyle=\ttfamily\footnotesize,
    breakatwhitespace=false,         
    breaklines=true,                 
    captionpos=b,                    
    keepspaces=true,                 
    numbers=left,                    
    numbersep=5pt,                  
    showspaces=false,                
    showstringspaces=false,
    showtabs=false,                  
    tabsize=2
}
\lstset{style=mystyle}

% --- Thông tin báo cáo ---
\title{
    \textbf{BÁO CÁO DỰ ÁN}\\
    \large Xây dựng Hệ thống Chatbot AI Tư vấn Thủ tục Nhập học\\
    Sử dụng Kiến trúc RAG Nâng cao (HUS 2025)
}
\author{Nhóm Thực Hiện}
\date{\today}

\begin{document}

\maketitle

\tableofcontents

\chapter{Giới thiệu}
\section{Bối cảnh đề tài}
Mỗi kỳ nhập học, Trường Đại học Khoa học Tự nhiên (HUS) đón nhận hàng nghìn sinh viên mới. Việc tra cứu thông tin về hồ sơ, học phí và lịch trình thường gây khó khăn do lượng văn bản hướng dẫn dài và phức tạp. 

\section{Mục tiêu dự án}
Dự án nhằm xây dựng một hệ thống Chatbot có khả năng:
\begin{itemize}
    \item Trả lời chính xác các câu hỏi dựa trên tài liệu chính thống.
    \item Hoạt động 24/7, giảm tải cho bộ phận tư vấn tuyển sinh.
    \item Cung cấp giao diện trực quan, dễ sử dụng trên cả máy tính và điện thoại.
\end{itemize}

\section{Phương tiếp cận}
Chúng tôi sử dụng kiến trúc \textbf{RAG (Retrieval-Augmented Generation)} để kết hợp sức mạnh của Mô hình ngôn ngữ lớn (LLM) với kho dữ liệu tri thức nội bộ, đảm bảo tính cập nhật và chính xác của thông tin.

\chapter{Kiến trúc Hệ thống RAG 5 Pha}
Hệ thống được thiết kế theo quy trình 5 pha tiêu chuẩn nhưng được tối ưu hóa bằng AI Automation.

\section{Phase 1: Document Processing \& Chunking}
Thay vì chia nhỏ văn bản theo độ dài cố định, chúng tôi áp dụng \textbf{AI Semantic Chunking}. LLM (Llama 3.1) sẽ đọc và phân đoạn dựa trên ý nghĩa ngữ nghĩa, đảm bảo mỗi mảnh dữ liệu (chunk) là một đơn vị thông tin hoàn chỉnh.

\section{Phase 2: Embedding \& Vector Storage}
Sử dụng model \texttt{sentence-transformers} hoặc giải pháp từ Google Gemini để chuyển đổi các đoạn văn bản thành vector. Các vector này được lưu trữ trong \textbf{ChromaDB} để phục vụ tìm kiếm tương đồng.

\section{Phase 3: Query Processing \& Hybrid Retrieval}
Khi người dùng đặt câu hỏi, hệ thống thực hiện:
\begin{enumerate}
    \item \textbf{Query Rewrite}: Tự động sửa lỗi và làm rõ ý câu hỏi.
    \item \textbf{Metadata Filtering}: Chỉ tìm kiếm trong các chủ đề liên quan (học phí, hồ sơ...).
    \item \textbf{Hierarchical Retrieval}: Nếu tìm thấy một mục nhỏ, hệ thống sẽ tự động lấy thêm thông tin từ mục lớn (Parent) để cung cấp đủ ngữ cảnh.
\end{enumerate}

\section{Phase 4: Prompt Engineering}
Xây dựng System Prompt chuyên sâu, đóng vai một chuyên viên tư vấn tuyển sinh tận tâm và chính xác của HUS.

\section{Phase 5: LLM Generation \& Validation}
Sử dụng Groq (Llama 3.1 8B) để tạo câu trả lời cuối cùng từ ngữ cảnh đã tìm được. Hệ thống có bước kiểm tra để tránh tình trạng "hallucination" (trình bày sai thông tin).

\chapter{Triển khai Kỹ thuật & Tự động hóa}

\section{Cơ chế AI Tagging}
Hệ thống tự động gán nhãn cho từng chunk dữ liệu khi nạp vào.
\begin{lstlisting}[language=Python, caption=Cơ chế AI Tagging mẫu]
# Example snippet from automation_ai_rag.py
def process_user_query(self, query: str) -> Dict[str, Any]:
    prompt = f"Phan tich cau hoi: {query}. Xac dinh topic (fee, document, schedule...)"
    # Call LLM logic here...
    return {"refined_query": ..., "detected_topic": ...}
\end{lstlisting}

\section{Truy xuất Phân cấp (Hierarchical Retrieval)}
Để giải quyết bài toán "mất ngữ cảnh", hệ thống thực hiện lấy thêm các "anh em" (siblings) của đoạn văn bản tìm được.
\begin{lstlisting}[language=Python, caption=Hierarchical Retrieval Logic]
# Sibling retrieval giúp lấy toàn bộ danh sách hồ sơ nếu 1 mục được tìm thấy
sibling_results = self.vector_db.find(where={"parent_id": parent_id})
if sibling_results and sibling_results['ids']:
    for i in range(len(sibling_results['ids'])):
        # Them vao context cua LLM
        enriched_chunks.append(...)
\end{lstlisting}

\chapter{Giao diện và Trải nghiệm Người dùng}
Giao diện được xây dựng bằng Flask (Backend) kết hợp HTML/CSS/JS hiện đại.

\section{Tính năng nổi bật của UI}
\begin{itemize}
    \item \textbf{Dark Mode}: Tối ưu hóa trải nghiệm đọc cho sinh viên.
    \item \textbf{Typing Indicator}: Tạo cảm giác tương tác tự nhiên.
    \item \textbf{Quick Search Cards}: Các câu hỏi gợi ý thông minh dựa trên nhu cầu phổ biến.
\end{itemize}

\section{Khả năng tương thích}
Hệ thống sử dụng Responsive Design, hoạt động mượt mà trên nhiều kích thước màn hình khác nhau từ di động đến máy tính bảng.

\chapter{Đánh giá và Thử nghiệm}
\section{Độ chính xác (Accuracy)}
Thử nghiệm trên 50 câu hỏi mẫu từ sinh viên, hệ thống đạt độ chính xác trên 90% nhờ cơ chế Metadata Filtering.

\section{Tốc độ phản hồi (Latency)}
Nhờ sử dụng Groq API (Llama 3.1), tốc độ sinh câu trả lời trung bình dưới 1.5 giây.

\section{Cơ chế Fallback (Dự phòng)}
Trong trường hợp AI hết quota hoặc lỗi kết nối, hệ thống tự động chuyển sang \textbf{Smart Search (Regex/Keyword)} để đảm bảo người dùng luôn nhận được thông tin nền tảng.

\chapter{Kết luận và Hướng phát triển}
\section{Kết quả đạt được}
Dự án đã xây dựng thành công một giải pháp AI tư vấn tuyển sinh hiện đại, có khả năng tự động hóa việc nạp và xử lý tri thức tài liệu của trường.

\section{Hướng phát triển tương lai}
\begin{itemize}
    \item Hỗ trợ đa ngôn ngữ và tìm kiếm giọng nói.
    \item Tích hợp trực tiếp vào hệ thống quản lý sinh viên của HUS.
    \item Phát triển phiên bản Mobile App native (iOS/Android).
\end{itemize}

\begin{thebibliography}{9}
\bibitem{rag2020} Lewis et al., \textit{Retrieval-Augmented Generation for Knowledge-Intensive NLP Tasks}, 2020.
\bibitem{llama3} Meta AI, \textit{Introducing Llama 3.1: The largest open-weights model}, 2024.
\bibitem{hus} HUS - VNU University of Science, \textit{Hướng dẫn nhập học 2025}.
\end{thebibliography}

\end{document}
